\section{Introduction}
\label{sec:introduction}

Aerial 3D scene understanding is central to autonomous uncrewed aerial vehicle (UAV) operation in cargo delivery, search and rescue, inspection, and human-UAV collaboration. In these settings, a perception system must infer not only visible image semantics, but also metric depth, free space, and the semantic occupancy of partially observed volumes. RGB cameras are attractive for UAVs because they are lightweight, inexpensive, and compatible with strict size, weight, and power constraints. However, monocular 3D perception remains difficult: image formation collapses scale, depth, and occlusion into a two-dimensional signal, and aerial views introduce altitude-dependent scale changes, oblique/top-down viewpoints, weak texture, and occlusion patterns that differ from ground-level driving scenes.

Most progress in semantic scene completion (SSC) and camera-based 3D perception has been driven by indoor and autonomous-driving benchmarks. SSCNet introduced semantic completion from depth input \cite{song2017sscnet}; MonoScene extended the task to monocular RGB \cite{cao2022monoscene}; and subsequent methods such as VoxFormer \cite{li2023voxformer}, TPVFormer \cite{huang2023tpvformer}, SurroundOcc \cite{wei2023surroundocc}, OpenOccupancy \cite{wang2023openoccupancy}, CGFormer \cite{yu2024cgformer}, and FoundationSSC \cite{chen2025foundationssc} improved view transformation, voxel attention, and semantic-geometric fusion. These designs demonstrate the value of explicit 3D representations.

OccuFly \cite{gross2026occufly} enables a focused study of this transfer problem. It provides aerial RGB images, metric depth maps, camera calibration, poses, altitude metadata, and semantic occupancy grids across multiple scenes and flight heights. Unlike 2D UAV segmentation datasets such as UAVid \cite{lyu2020uavid} or reconstruction-oriented resources such as UrbanScene3D \cite{lin2022urbanscene3d}, OccuFly supports joint evaluation of monocular depth estimation (MDE), multiview geometry, and SSC from real aerial observations. We use it to ask whether modern foundation representations provide transferable geometric and semantic priors for UAV 3D perception, and whether calibrated neighboring views improve SSC beyond monocular depth priors.

We investigate frozen V-JEPA~2.1 \cite{murlabadia2026vjepa21} image and video features for aerial MDE and depth-aware SSC. V-JEPA-style models learn by predicting latent representations rather than reconstructing pixels \cite{bardes2024vjepa,assran2025vjepa2}; V-JEPA~2.1 further emphasizes dense predictive supervision, deep self-supervision, and unified image/video tokenization. Rather than finetuning the backbone, we keep it frozen and vary only the decoder, metadata conditioning, temporal adapter, and lifting mechanism. This protocol makes the study diagnostic: improvements can be attributed to the downstream structure used to extract depth and occupancy from the representation.

Our evaluation has three parts. First, we compare linear probes, DPT-style multiscale decoders \cite{ranftl2021dpt}, Simple Feature Pyramid (SFP) adapters inspired by ViTDet \cite{li2022vitdet}, and Feature-wise Linear Modulation (FiLM) \cite{perez2018film} using altitude and camera intrinsics. This separates linear depth accessibility, dense reconstruction capacity, final-layer pyramid adaptation, and metric scale conditioning. Second, we compare RGB features with spatio-temporal V-JEPA~2.1 video features using a dense temporal-attention adapter that converts tubelet tokens into target-frame maps. Third, we evaluate SSC variants inspired by CGFormer and FoundationSSC, replacing or augmenting their geometry pathways with V-JEPA-derived depth probabilities and a pose-aware multiview stereo (MVS) cost-volume branch based on calibrated UAV views.

Our contributions are:
\begin{itemize}
    \item A systematic frozen-backbone evaluation of V-JEPA~2.1 RGB and video representations for aerial metric depth estimation on OccuFly.
    \item A controlled comparison of linear probes, DPT-style decoders, SFP-adapted final-layer decoders, and altitude/intrinsics-conditioned FiLM variants.
    \item An analysis of spatio-temporal feature transfer through dense temporal attention over V-JEPA~2.1 video tokens.
    \item A depth-to-SSC study that injects V-JEPA-derived depth maps and depth distributions into CGFormer- and FoundationSSC-inspired aerial SSC variants.
    \item A pose-aware FoundationSSC variant that replaces monocular depth lifting with calibrated multiview cost-volume depth estimation from neighboring UAV frames.
\end{itemize}
